\documentclass[10pt,a4paper]{amsart}
\usepackage[margin=19mm]{geometry}
\usepackage{amsmath,amssymb,mathtools}
\usepackage[scr=rsfs,cal=euler]{mathalfa}
\usepackage{xfrac}
\usepackage[unicode,hidelinks,bookmarks=false]{hyperref}
\newcommand{\ie}{i.e.\ }
\newcommand{\cf}{cf.\ }
\newcommand{\resp}{respectively\ }
\newcommand{\Subsection}{\subsection}
\providecommand{\nobreakdash}{\nobreak\hbox{-}}
\setlength{\emergencystretch}{2em}
\multlinegap0pt
\usepackage{amssymb}
\newtheorem{prop}{Proposition}
\newtheorem{lemma}{Lemma}
\title{Large fluctuations of sums of a random multiplicative function}
\author{Besfort Shala}
\date{Source: arXiv:2602.20086v3 (CC BY 4.0)}
\begin{document}
\maketitle

\noindent\emph{Source and licence.} Large fluctuations of sums of a random multiplicative function. Version arXiv:2602.20086v3 (CC BY 4.0), distributed by arXiv under the Creative Commons Attribution 4.0 International licence, http://creativecommons.org/licenses/by/4.0/, as recorded in the arXiv licence metadata for this exact version and shown on the abstract page https://arxiv.org/abs/2602.20086v3. Full licence notice: \texttt{licenses/arxiv-2602.20086-CC-BY-4.0.txt}.\par\smallskip

\noindent\textit{Editorial scope.} Counted unit: the article's lemma lemma:tooManyPrimes (source0.tex lines 279--284). The article's complete original proof of it is delivered with it (source0.tex lines 286--296, one \texttt{proof} environment of 11 lines). No mathematics is written, repaired or re-derived: the statement and the proof are the article's own lines, in the article's own notation, and no retained line is substituted or re-flowed.  Retained uncounted as the source-local context the proof needs: lines 254--260.  Nothing was omitted from any retained span, so the package carries no omission record.  No retained line contains a citation, so the wrapper carries no bibliography and no bibliography entry.  No retained line contains a figure environment, an image-inclusion command or drawing code, so no image is delivered and the package declares no asset dependency.  Editorial formatting only, in addition: an AMS \texttt{amsart} preamble that restates the article's own declarations and macros used by the retained lines, namely the environments \texttt{prop}, \texttt{lemma} on separate counters, together with the standard packages those lines need.  The retained environments print in this document as Proposition 1, Lemma 1 in order of appearance, whereas the article prints the same items with the numbering of its own full document.  The complete native source of the article as distributed in the arXiv e-print for this exact version is bundled unchanged as \texttt{sources/195157/source0.tex}.
\begin{prop}\label{prop:nair-tenenbaum}
    Let $F$ be a non-negative multiplicative function such that there exists a constant $A\geq 1$ and $\epsilon >0$ such that $F(m)\leq Am^{\epsilon}$ for all $m\in\mathbb N$. Let $Q\in\mathbb Z[x]$ be a fixed polynomial of degree $d$ with no fixed prime factor, and let $\rho(n)$ be the number of solutions to $Q(x)\equiv 0\pmod n$. Suppose that $Q$ has $r$ distinct irreducible factors $Q_1, Q_2, \ldots, Q_r$, and write $Q = \prod_{i=1}^r Q_i^{\gamma_i}$. Let $\rho_i(n)$ be the number of solutions to $Q_i(x)\equiv 0\pmod n$ for $1\leq i\leq r$. Denote the discriminant of the square-free kernel of $Q$ by $D$. If $\epsilon<\frac{1}{8d^2}$, then 
    \begin{multline*}
        \sum_{N<n\leq N+H} F(\lvert Q(n)\rvert ) \\ \ll H\prod_{p\leq N}\left(1-\frac{\rho(p)}{p}\right) \sum_{\substack{n_1^{\gamma_1}n_2^{\gamma_2}\cdots n_r^{\gamma_r}\leq N \\ (n_i, n_j)=1 \text{ }\forall i\neq j \\ (n_i, D)=1 \text{ }\forall i}} F(n_1^{\gamma_1}\cdots n_r^{\gamma_r})\frac{\rho_1(n_1)\rho_2(n_2)\cdots \rho_r(n_r)}{n_1n_2\cdots n_r},
    \end{multline*}
    whenever $N$ is large enough, and $N^{4d^2\epsilon}\leq H\leq N$. 
\end{prop}

\begin{lemma}\label{lemma:tooManyPrimes}
    Let $\varepsilon>0$ and let $N$ and $H$ be large enough such that $H\geq N^{5\varepsilon}$. Let $\Omega(n)$ denote the number of prime factors of $n$, counted with multiplicity. Then there exists $\varepsilon'>0$ such that 
    \begin{equation*}
        \#\{N-H<n\leq N : \Omega(n)>(1+\varepsilon)\log\log N\}\ll \frac{H}{(\log N)^{\varepsilon'}}.
    \end{equation*}
\end{lemma}

\begin{proof}
The quantity we want to bound is certainly at most 
\begin{equation*}
    \frac{1}{\exp((1+\varepsilon)\log\log N\log(1+\varepsilon))}\sum_{N-H<n\leq N} (1+\varepsilon)^{\Omega(n)}.
\end{equation*}
Applying Proposition \ref{prop:nair-tenenbaum} with  $A=1, \epsilon = \log(1+\varepsilon)/\log 2=\varepsilon/\log 2 +O(\varepsilon^2)$ to the sum gives 
\begin{equation*}
    \sum_{N-H<n\leq N} (1+\varepsilon)^{\Omega(n)}\ll H(\log N)^{\varepsilon} = H\exp(\varepsilon\log\log N),
\end{equation*}
whenever $H\geq N^{4\epsilon}\gg N^{5\varepsilon}$. Putting the two bounds together gives that the quantity we want to bound is $\ll H/(\log N)^{(1+\varepsilon)\log(1+\varepsilon)-\varepsilon}$, so we may take $\varepsilon' = (1+\varepsilon)\log(1+\varepsilon)-\varepsilon = \frac{\varepsilon^2}{2} + O(\varepsilon^3)$. 
\end{proof}

\end{document}
